\begin{abstract}
We give a minimal algebraic account of evolution and extension in an
ordered encounter. The visible encounter algebra is the Klein four
group \(E=\{1,a,b,ab\}\), where \(a\) and \(b\) register two
participating roles. Because \(E\) is commutative, it records
participation but not encounter order. We lift \(a\) and \(b\) to
generally noncommuting elements \(A\) and \(B\). The products \(AB\)
and \(BA\) have the same visible image, while
\[
[A,B]=ABA^{-1}B^{-1}
\]
belongs to the quotient kernel and satisfies \(AB=[A,B]BA\). We call the
quotient-visible product evolution and the retained kernel residue
extension.

When \(A\) and \(B\) are involutions, \(H=AB\) has inverse \(BA\) and
\([A,B]=H^2\). We exhibit the finite realization
\[
\langle A,B\mid A^2=B^2=1,(AB)^8=1\rangle\cong D_{16}.
\]
Its quotient onto \(E\) has kernel \(\langle H^2\rangle\) of order
four, while the \(H\)-orbit is \(C_8\). A further \(C_8\)-to-\(C_4\)
quotient returns visibly at \(H^4\) while
retaining a nontrivial lifted state. The result is algebraic:
extension means retained kernel information, not spatial expansion or
physical backaction. No physical theory is claimed.
\end{abstract}
